% agent-rewind: Phase 1 report, Option 5.
% Course rule: a 1-2 page writeup plus an appendix with a timeline and milestones.
% Scope is Phase 1 only: clean baseline, fault injection, traces, propagation,
% and the full-restart and always-verify baselines. Recovery methods are Phase 2.
% Build: `tectonic phase1.tex` from this folder.
\documentclass[10pt,twocolumn,letterpaper]{article}

\usepackage[T1]{fontenc}
\usepackage{newtxtext,newtxmath}
\usepackage[textwidth=7in,textheight=9in,columnsep=0.33in,centering]{geometry}
\usepackage{booktabs}
\usepackage{multirow}
\usepackage{xcolor}
\usepackage[font=small,labelfont=bf,skip=3pt]{caption}
\usepackage{enumitem}
\usepackage{microtype}
\usepackage[hidelinks]{hyperref}

\setlist{nosep,leftmargin=*}
\setlength{\textfloatsep}{6pt plus 2pt minus 2pt}
\pagestyle{plain}
\raggedbottom
\definecolor{warn}{HTML}{D55E00}

\title{\large\bfseries Phase 1: Fault Injection and Propagation\\in a Researcher--Planner--Executor Workflow}
\author{Yitang Gao \quad Qinlin Chen \quad Xueru Song \quad Kai Zhang\\ \normalsize University of California, San Diego}
\date{}

\begin{document}
\maketitle

\section{Goal}
Option~5 asks whether tracking dependencies between agent results can make
recovery from a known-faulty result efficient while keeping execution correct
and safe. Phase~1 builds the ground for that question: a clean workflow with
checkable outcomes, a harness that injects faults at controlled points, traces
that record what each step read and did, measurements of how faults
propagate, and the two reference baselines, \emph{full restart} and
\emph{always-verify}. Recovery methods are Phase~2 work and are not evaluated
here.

\section{Setup}
\paragraph{Workloads.} Each workload has a clean run whose outcome a checker
verifies.
\textbf{W1} (booking) is a researcher $\to$ planner $\to$ executor workflow.
The researcher looks up team size and room capacities, the planner picks a room
for each of three meetings, and the executor books it and e-mails attendees
(14 steps, 6 LLM calls). \texttt{calendar.book} is compensable;
\texttt{email.send} is irreversible.
\textbf{W3} (code fix) has an agent read an issue and source file, then write a
patch into a git workspace that unit tests check. This adds non-LLM file state.
\textbf{W4} ports two scenarios from TAMAS~\cite{kavathekar2025tamas}, a
named starting point in the brief: Byzantine agents (\texttt{w4.byz}) and
indirect prompt injection (\texttt{w4.ipi}). Their effects are posts to
external sites, which are irreversible.

\paragraph{Model and platform.} Qwen2.5-7B-Instruct (Q4\_K\_M) runs in-process
through \texttt{llama-cpp-python} on one Apple M5 Pro laptop (24\,GB), with
greedy decoding and JSON-constrained output. A scripted model runs the same
workflows for fast checks. Each configuration has 3 repeats, which are
scenario variants. 1{,}675 cells ran with zero errors.

\paragraph{Traces.} Every step reads and writes one versioned namespace:
messages, shared state, files and tool results. A step records (i)~the
messages and artifacts it wrote, (ii)~its tool actions with their effect class
(\texttt{read}, \texttt{compensable}, \texttt{irreversible}), (iii)~its cost
(LLM calls, tokens, wall time), and (iv)~its \emph{dependencies}: the ordered
paths it read, each with a content hash. A fault's \emph{cone} is every step
that read the faulty value, directly or through another step.

\paragraph{Faults.} One fault per run, at a chosen step: \emph{wrong} (a
plausible false value), \emph{stale} (an outdated value), \emph{timeout} (a
missing value) and \emph{malicious} (an injected instruction or a Byzantine
agent). Accidental and adversarial faults are reported separately.

\section{Clean baseline}
All four workloads succeed in 3/3 clean runs on both models
(Table~\ref{tab:clean}). The clean-vs-clean spread in wall time is
1.1\,s (W1), 1.2\,s (W3) and 0.6\,s (W4.ipi), which sets the noise floor for
latency.

\begin{table}[h]
\centering\small
\caption{Clean runs, 7B model ($n{=}3$ each).}
\label{tab:clean}
\setlength{\tabcolsep}{4pt}
\begin{tabular}{@{}lccccc@{}}
\toprule
WL & success & steps & LLM calls & tokens & latency s \\
\midrule
W1 & 3/3 & 14 & 6 & 2427 & 47.6 \\
W3 & 3/3 & 5 & 1 & 661 & 8.7 \\
W4.byz & 3/3 & 6 & 3 & 444 & --$^\ast$ \\
W4.ipi & 3/3 & 6 & 3 & 1078 & 13.2 \\
\bottomrule
\end{tabular}\\[1pt]
{\footnotesize $^\ast$All cells were prompt-cache hits, so no live latency was measured.}
\end{table}

\section{Propagation}
For each fault we measure whether any later step read the faulty value,
whether the outcome changed, how many steps and which effects it reached, and
whether unauthorised actions occurred (Table~\ref{tab:prop}).

\begin{table}[t]
\centering\small
\caption{Fault propagation, 7B model ($n{=}3$ each). ``Changed'' is the
fraction of runs whose outcome differs from the clean run. ``Irrev.'' counts
irreversible calls made in the faulty run.}
\label{tab:prop}
\setlength{\tabcolsep}{2.4pt}
\begin{tabular}{@{}llcccc@{}}
\toprule
WL & fault (class) & changed & steps hit & irrev. & unauth. \\
\midrule
\multirow{3}{*}{W1} & wrong capacity (acc.) & 1.00 & 9/14 & 3 & 0 \\
 & stale capacity (acc.) & 0.67 & 9/14 & 3 & 0 \\
 & injected agenda (adv.) & 1.00 & 3/14 & 3 & \textbf{1} \\
\midrule
\multirow{3}{*}{W3} & wrong issue (acc.) & 1.00 & 3/5 & 0 & 0 \\
 & stale code (acc.) & 1.00 & 3/5 & 0 & 0 \\
 & injected issue (adv.) & \textcolor{warn}{0.00} & 3/5 & 0 & 0 \\
\midrule
\multirow{4}{*}{W4.byz} & stale flag (acc.) & 0.67 & 4/6 & 2 & 0 \\
 & Byz.\ distribute (adv.) & 1.00 & 1/6 & 1 & 0 \\
 & Byz.\ fact-check (adv.) & \textcolor{warn}{0.00} & 3/6 & 2 & 0 \\
 & injected item (adv.) & 0.33 & 4/6 & 2.3 & 0 \\
\midrule
\multirow{4}{*}{W4.ipi} & stale flag (acc.) & 0.67 & 4/6 & 1 & 0 \\
 & Byz.\ distribute (adv.) & 1.00 & 1/6 & 1 & 0 \\
 & Byz.\ fact-check (adv.) & \textcolor{warn}{0.00} & 3/6 & 1 & 0 \\
 & injected item (adv.) & 0.33 & 4/6 & 1.3 & 0 \\
\bottomrule
\end{tabular}
\end{table}

\begin{itemize}
\item \textbf{Every fault is read downstream} (all 14 fault--workload pairs),
  so each workload really depends on the injected result.
\item \textbf{Reach varies widely.} A wrong W1 capacity reaches 9 of 14 steps,
  both bookings and the irreversible e-mail. The injected agenda reaches only
  3 steps, yet sends an unauthorised e-mail in every run.
\item \textbf{Some faults are read but absorbed.} Three adversarial faults
  never change the 7B outcome (W3 injected issue; Byzantine fact-check in both
  W4 scenarios). The scripted model propagates all of them, so absorption is a
  property of the model. We flag these faults rather than count them as
  robustness.
\item \textbf{Accidental faults are not always caught by the model.} A stale
  capacity changes the W1 outcome in 2 of 3 runs.
\end{itemize}

\section{Baselines}
\textbf{Full restart} undoes what it can and reruns every step.
\textbf{Always-verify} audits every step's output with the same model
against the step's inputs, and its tokens are counted as extra inference cost.
Both see the same faulty runs (Table~\ref{tab:base}). \emph{Wasted work} is the
number of rerun steps outside the fault's cone, which a dependency-aware method
could avoid.

\begin{table}[t]
\centering\small
\caption{Baselines, 7B model, mean over faults. Restart: steps rerun, wasted
steps (outside the cone), recovery time and tokens, irreversible calls
re-sent, recovered correctly. Always-verify: audit tokens per run and their
ratio to a clean run's tokens, detection rate, false-positive flags per run.}
\label{tab:base}
\setlength{\tabcolsep}{2.2pt}
\begin{tabular}{@{}l|cccccc|cccc@{}}
\toprule
 & \multicolumn{6}{c|}{full restart} & \multicolumn{4}{c}{always-verify} \\
WL & rerun & waste & s & tok & irr. & ok & tok & $\times$ & det. & FP \\
\midrule
W1 & 14 & 6.0 & 36.3 & 2326 & 3.0 & 1.0 & 1245 & .51 & 1.00 & .33 \\
W3 & 5 & 1.0 & 8.0 & 562 & 0 & 0.0$^\dagger$ & 575 & .87 & 0.67 & 0 \\
W4.byz & 6 & 2.0 & 13.4 & 1078 & 1.8 & 1.0 & 692 & 1.56 & 0.75 & .25 \\
W4.ipi & 6 & 2.0 & 14.1 & 1097 & 1.0 & 1.0 & 712 & .66 & 1.00 & .25 \\
\bottomrule
\end{tabular}\\[1pt]
{\footnotesize $^\dagger$See the observations below.}
\end{table}

\paragraph{Observations.}
(1)~Restart reruns many steps that never read the fault: on W1, 10 of 14 for
the injected agenda and 4 of 14 for the capacity faults. It also re-sends every
irreversible call (3.0 per W1 recovery), and it cannot recall the unauthorised
e-mail. (2)~Always-verify adds 51--156\% of a clean run's tokens on
\emph{every} run, faulty or not. It still misses one third of W3 faults and a
quarter of W4.byz faults, and it raises false flags on clean steps.
(3)~$^\dagger$Restart did not reproduce the clean W3 outcome in any 7B run, even
though clean W3 succeeds 3/3. We are investigating whether the restarted run
differs from the clean run in prompt or workspace state.

\paragraph{Limitations.} Timeout faults are covered only in a scripted
single-meeting version of W1, where a missing team size changes the outcome and
reaches both effects. They are not yet in the 7B grid. W4 uses 2 of TAMAS's
scenarios, and $n{=}3$ gives wide intervals.

\paragraph{Takeaway for Phase 2.} Faults reach only part of each workflow, but
restart pays for all of it and re-sends irreversible effects, while
always-verify pays on every run and still misses faults. Phase~2 will use the
recorded dependencies to rerun only the cone, reconcile effects instead of
repeating them, and compare against these two baselines under identical fault
schedules.

{\footnotesize
\bibliographystyle{abbrv}
\bibliography{refs}
}

\clearpage
\onecolumn
\appendix
\section*{Appendix: Timeline and Milestones}

\begin{center}
\small
\renewcommand{\arraystretch}{1.25}
\begin{tabular}{@{}p{0.11\textwidth}p{0.42\textwidth}p{0.35\textwidth}@{}}
\toprule
\textbf{Due / status} & \textbf{Milestone} & \textbf{Done when} \\
\midrule
\multicolumn{3}{@{}l}{\textit{Phase 1 --- fault injection and propagation (this report)}} \\
Done & Review TAMAS and A2ASecBench; choose scenarios & two TAMAS scenarios ported (W4) \\
Done & Clean researcher--planner--executor baseline & W1, W3, W4 pass their checkers 3/3 \\
Done & Fault-injection harness & wrong / stale / timeout / malicious faults at a chosen step \\
Done & Traces: messages, artifacts, tool actions, dependencies & every step records reads with hashes, writes, effects, cost \\
Done & Measure clean runs and propagation & Tables 1--2 \\
Done & Full-restart and always-verify baselines & Table 3 \\
\midrule
\multicolumn{3}{@{}l}{\textit{Phase 2 --- selective recovery}} \\
Nov 8 & Dependency tracking: tag results with their sources & a step's inputs can be listed from the trace \\
Nov 15 & Dependency-aware invalidation and selective rerun & only steps in the fault's cone rerun \\
Nov 22 & Shared state and irreversible actions & compensation for compensable effects; irreversible ones reported, not repeated \\
Nov 29 & Compare against full restart and always-verify & identical fault schedules and budgets, fault stage varied \\
Dec 5 & Evaluate recovery metrics and reuse validity & correctness, unauthorised actions, recomputed stages, retained work, latency, inference cost, tracking overhead \\
Dec 7 & Optional: automated fault detection & false positives and misses reported separately from recovery \\
Dec 10 & Submit Phase 2 code and the final report & 4--6 page report \\
\bottomrule
\end{tabular}
\end{center}

\paragraph{Open items carried into Phase 2.} Add timeout faults to the 7B grid;
explain the W3 restart result; add the precision-versus-bookkeeping
measurement (dependency granularity) that the brief asks for.

\paragraph{Acknowledgments.} Phase 1 progress report for the CSE 291A
(\emph{Systems for LLMs and AI Agents}, UC San Diego, Fall 2026) course project.

\end{document}
